\chapter{Cronograma de atividades}
\label{cap:cronograma}

\section{Períodos de referência}

Adoto como referência o calendário de 2026 que recebi para organizar o TCC I. Na Tabela~\ref{tab:calendario}, reproduzo as datas escritas nesse material. Não considero uma etapa concluída apenas porque seu período transcorreu; registrarei a realização de cada entrega.

\begin{table}[htb]
\centering
\caption{Períodos indicados no calendário fornecido}
\label{tab:calendario}
\small
\begin{tabular}{p{0.24\textwidth}p{0.32\textwidth}p{0.33\textwidth}}
\hline
\textbf{Período de 2026} & \textbf{Etapa} & \textbf{Entrega associada} \\ \hline
10/08 a 22/08 & Definição de tema e orientador. & Tema delimitado e orientação definida. \\ \hline
24/08 a 26/09 & Escrita dos trabalhos de TCC I e II. & Versão escrita correspondente à etapa acadêmica do estudante. \\ \hline
26/09 a 07/11 & Entrega do trabalho à banca. & Arquivo revisado e registro do encaminhamento. \\ \hline
03/11 a 28/11 & Defesas. & Apresentação, avaliação e registro das recomendações. \\ \hline
30/11 a 04/12 & Correções e envio da versão final no SIGAA. & Versão corrigida e comprovante de envio. \\ \hline
\end{tabular}
\par\smallskip
{\footnotesize Fonte: calendário acadêmico fornecido ao autor; transcrição das datas textuais.}
\end{table}

O material apresenta divergência entre as datas escritas, as durações aproximadas e a marcação gráfica: o encerramento da escrita e o início da entrega aparecem como 26/09, embora a marcação se estenda por outubro. Até confirmar esse ponto, preservo a data textual de 26/09 como marco conservador para preparar a primeira versão. Confirmarei com o orientador a data individual de defesa e a antecedência exigida para envio à banca; não interpreto o limite de 07/11 como autorização para encaminhar o trabalho depois da defesa.

\section{Entregáveis do projeto de TCC I}

Como este documento corresponde ao projeto de TCC I, concentro as entregas desta etapa na delimitação da pesquisa, na fundamentação, na descrição do estado inicial do protótipo e no protocolo de avaliação. Na Tabela~\ref{tab:cronograma}, estabeleço marcos internos compatíveis com a preparação de uma primeira versão até 26/09. Essas datas representam meu planejamento, e não novos prazos institucionais.

\begin{table}[htb]
\centering
\caption{Marcos internos propostos para o projeto de TCC I}
\label{tab:cronograma}
\small
\begin{tabular}{p{0.19\textwidth}p{0.38\textwidth}p{0.32\textwidth}}
\hline
\textbf{Prazo em 2026} & \textbf{Entrega verificável} & \textbf{Critério de conclusão} \\ \hline
Até 18/09 & Introdução fundamentada e levantamento de sistemas existentes. & Fontes citadas e escopo coerente com o problema. \\ \hline
Até 22/09 & Requisitos iniciais e quadro de estado do protótipo. & Separação entre código já desenvolvido, correções e avaliação futura. \\ \hline
Até 24/09 & Diagramas de arquitetura e dados e protocolo de testes. & Correspondência com o código e cenários com resultados esperados. \\ \hline
Até 26/09 & Primeira versão integral do projeto para revisão. & Objetivos, método, cronograma e referências articulados. \\ \hline
26/09 a 30/10 & Revisões orientadas e versão candidata à banca. & Pendências da revisão tratadas e PDF conferido. \\ \hline
Até 07/11 & Encaminhamento à banca, antecipado quando necessário. & Registro de envio antes da defesa e na antecedência confirmada. \\ \hline
03/11 a 28/11 & Defesa na data individual designada. & Apresentação e lista de correções solicitadas. \\ \hline
30/11 a 04/12 & Revisão final e envio no SIGAA. & Correções incorporadas e envio registrado. \\ \hline
\end{tabular}
\par\smallskip
{\footnotesize Fonte: elaboração própria, tomando o calendário fornecido como referência.}
\end{table}

\section{Execução técnica e dependências}

Executarei o desenvolvimento complementar e a coleta de resultados no TCC II. Como ainda preciso compatibilizar as datas dessa etapa com a orientação, apresento na Tabela~\ref{tab:execucao-tecnica} uma sequência de execução, sem atribuir a ela prazos institucionais ainda não confirmados. No TCC I, trato esses itens apenas como planejamento e não os apresento como resultados concluídos.

\begin{table}[htb]
\centering
\caption{Sequência proposta para desenvolvimento e avaliação no TCC II}
\label{tab:execucao-tecnica}
\small
\begin{tabular}{p{0.21\textwidth}p{0.40\textwidth}p{0.28\textwidth}}
\hline
\textbf{Etapa do TCC II} & \textbf{Atividade e entrega} & \textbf{Dependência} \\ \hline
Etapa 1 & Consolidar ambiente, dados sintéticos e correções do pipeline. Entregar versão de trabalho reproduzível. & Requisitos e protocolo definidos. \\ \hline
Etapa 2 & Executar testes preliminares e corrigir falhas demonstradas. Entregar versão candidata à avaliação. & Ambiente e testes disponíveis. \\ \hline
Etapa 3 & Executar avaliação funcional e de qualidade. Entregar dados brutos e relatório de testes. & Versão identificada e critérios fixados. \\ \hline
Etapa 4 & Analisar resultados e limitações. Entregar síntese e documentação da versão avaliada. & Evidências de execução preservadas. \\ \hline
\end{tabular}
\par\smallskip
{\footnotesize Fonte: elaboração própria; sequência técnica sujeita ao calendário do TCC II e à orientação.}
\end{table}

Quando o calendário do TCC II estiver confirmado, atribuirei datas a essas etapas preservando a sequência de dependências. Na defesa do projeto de TCC I, não apresentarei como concluídos resultados ainda não coletados. Para a defesa do TCC II, concluirei a avaliação e a análise antes do encaminhamento à banca; se houver atraso na versão do protótipo, revisarei explicitamente o escopo e as datas, sem omitir testes pendentes.

Serei responsável pela produção dos artefatos, com acompanhamento do orientador na delimitação e na revisão das entregas. Em cada marco, registrarei data prevista, data realizada, versão ou documento entregue e pendências. Assim, utilizarei o cronograma para acompanhar produtos concretos e alterações justificadas no planejamento.
